\documentclass[conference]{IEEEtran}

\usepackage{amsmath,amssymb}
\usepackage{newtxtext,newtxmath}
\usepackage{graphicx}
\usepackage{booktabs}
\usepackage{array}
\usepackage{cite}
\usepackage{url}
\usepackage[hidelinks]{hyperref}
\usepackage[caption=false,font=footnotesize]{subfig}
\usepackage{orcidlink}

\begin{document}

% -------------------------------------------------------
% TITLE
% -------------------------------------------------------

\title{
\textbf{\textit{StratoWatch: Air Quality Forecasting}}\\[3pt]
\textbf{\textit{with Spatio-Temporal Transformers and Graph-Aware Modeling}}
}

\author{
\IEEEauthorblockN{1\textsuperscript{st} Siddharth Bose \orcidlink{0009-0009-5156-5996}}
\IEEEauthorblockA{
School of Computer Science and Technology\\
Bennett University, Greater Noida, India\\
\href{mailto:s24cseu0543@bennett.edu.in}{s24cseu0543@bennett.edu.in}}
}

\maketitle

% -------------------------------------------------------
% ABSTRACT
% -------------------------------------------------------

\begin{abstract}
Forecasting ground-level ozone (O$_3$) and nitrogen dioxide (NO$_2$) a few hours ahead is
difficult because pollutant concentrations depend on both recent atmospheric conditions and
rapidly changing meteorological patterns. StratoWatch is a forecasting and evaluation
framework motivated by ISRO Problem Statement 25178 that studies this problem at both
single-site and multi-site scales. At the single-site level, we compare a forecast-only
baseline with Random Forest, XGBoost, LSTM, TCN, and a temporal Transformer using a 24-hour
input window and a 6-hour forecast horizon. XGBoost gives the strongest single-site result,
with an MAE of 0.4039 and an $R^2$ of 0.3848 in the near-normalized evaluation units. We then
extend the problem to seven Delhi monitoring stations using a spatio-temporal Transformer
and two graph-based variants. The plain ST-Transformer obtains the lowest MAE at
18.8349~$\mu$g/m$^3$, while the Dynamic Wind Graph ST-Transformer obtains the lowest RMSE
of 26.9455~$\mu$g/m$^3$ and the highest $R^2$ of 0.3911. These results suggest that explicit
spatial modeling is useful, although its benefit depends on the evaluation metric rather than
producing a uniform improvement. The current system is a short-horizon research prototype
rather than a complete 24--48 hour operational forecasting system. Alongside the experiments,
StratoWatch provides a React/FastAPI interface for running the forecasting pipelines and
viewing their results.
\end{abstract}

\begin{IEEEkeywords}
air quality forecasting, ozone, nitrogen dioxide, transformers,
spatio-temporal forecasting, graph neural networks, Delhi, ISRO
\end{IEEEkeywords}

% -------------------------------------------------------
\section{Introduction}
% -------------------------------------------------------

Air-quality forecasting is difficult because pollutant concentrations are affected by several
processes at the same time. O$_3$ and NO$_2$ respond to recent concentration history, wind,
temperature, humidity, boundary-layer conditions, and other atmospheric variables. These
effects can change considerably from one hour to the next. This makes short-term forecasting
different from simply predicting a smooth time series: a model also has to deal with sudden
changes and pollution peaks.

Physics-based forecasts provide an important starting point for this problem, but their
predictions can still contain systematic errors at individual locations. Machine learning
offers a way to learn patterns in historical observations and environmental variables that may
not be fully captured by the initial forecast.

This work was motivated by ISRO Problem Statement 25178, which focuses on automated forecasting
of surface-level O$_3$ and NO$_2$ using artificial intelligence and machine learning. The
implementation presented here should be viewed as a research prototype motivated by that
problem rather than a complete implementation of its 24--48-hour operational requirement.
The experiments currently use a six-hour forecast horizon so that the different model
architectures can be compared under the same evaluation setting.

The study has two main parts. The first evaluates forecasting models at a single-site level.
The second moves to a seven-site setting and asks whether explicitly representing relationships
between monitoring stations improves forecasting performance.

The second question is particularly interesting because nearby stations are not independent.
Pollution can move between locations, and wind can provide information about the direction of
that movement. At the same time, meteorological variables already provide some of this
information directly to a forecasting model. It is therefore not obvious how much additional
information an explicit spatial graph will provide.

The main contributions of this work are:

\begin{enumerate}
\item A single-site forecasting pipeline comparing tree-based, recurrent, convolutional,
and Transformer models for six-hour O$_3$ and NO$_2$ forecasting.
\item A seven-station spatio-temporal forecasting pipeline using a shared Transformer
representation across monitoring stations.
\item Static and wind-aware dynamic graph variants for testing the effect of explicit
spatial relationships between stations.
\item A reproducible evaluation structure based on frozen model artifacts and consolidated
evaluation results.
\item A working web-based interface that connects the forecasting models to a FastAPI
backend and allows results to be inspected without running the research code manually.
\end{enumerate}

% -------------------------------------------------------
\section{Related Work}
% -------------------------------------------------------

\subsection{Machine Learning for Air Quality Forecasting}

Machine learning has been widely used for air-quality forecasting because pollutant
concentrations depend on nonlinear relationships between historical measurements and
environmental variables. Tree-based methods are particularly useful when the input contains
many engineered variables and nonlinear interactions. XGBoost provides an efficient gradient
boosting framework for this type of tabular learning problem \cite{xgboost}.

\subsection{Sequence Modeling}

Recurrent models such as Long Short-Term Memory (LSTM) networks are designed to capture
dependencies in sequential data and have therefore been used extensively for environmental
forecasting \cite{lstm}. Temporal Convolutional Networks (TCNs) provide another approach,
using causal and dilated convolutions to model dependencies over different temporal scales
\cite{tcn}.

\subsection{Transformers for Time-Series Forecasting}

Transformers use self-attention to allow information from different positions in a sequence to
interact directly \cite{vaswani}. This makes them attractive for forecasting problems where
recent and older observations may both contain useful information. Informer further explored
efficient Transformer-based forecasting for long sequences \cite{informer}. In StratoWatch,
Transformer encoders are used to model temporal dependencies over a fixed 24-hour input
window.

\subsection{Spatio-Temporal Graph Models}

When observations come from multiple monitoring locations, the spatial relationship between
stations becomes another possible source of information. Graph-based forecasting methods
represent locations as nodes and use edges to exchange information between them. Examples
include Graph WaveNet, AGCRN, and DCRNN \cite{graphwavenet,agcrn,dcrnn}. Recent air-quality
work has also explored combinations of graph and Transformer architectures, including
CEEMDAN-GNN-Transformer and extreme-event-oriented spatio-temporal graph models
\cite{ceemdan2025,estgcn2025}.

These approaches motivate the multi-site experiments in StratoWatch, where spatial graph
information is compared against a Transformer that already receives meteorological and
multi-station information.

% -------------------------------------------------------
\section{Problem Formulation}
% -------------------------------------------------------

For the single-site experiment, the model receives a historical window of observations and
engineered environmental features. Let

\begin{equation}
X_t \in \mathbb{R}^{T \times F}
\end{equation}

denote the input sequence, where $T=24$ hours and $F=134$ input features. The model predicts
both pollutants over the next six hours:

\begin{equation}
\hat{Y}_t \in \mathbb{R}^{H \times C},
\end{equation}

where $H=6$ and $C=2$, corresponding to O$_3$ and NO$_2$.

The final single-site XGBoost model uses direct target prediction rather than residual
correction. The 24-hour history is flattened into a 3,216-dimensional feature vector
($24\times134$), and twelve independent regressors are used for the six forecast horizons and
two target pollutants.

For the multi-site experiment, the input is represented as

\begin{equation}
X \in \mathbb{R}^{T \times S \times F},
\end{equation}

where $T=24$ hours, $S=7$ monitoring stations, and $F=28$ features per station. The output is

\begin{equation}
\hat{Y} \in \mathbb{R}^{H \times S \times C},
\end{equation}

where $H=6$ forecast hours and $C=2$ target pollutants.

% -------------------------------------------------------
\section{Dataset and Preprocessing}
% -------------------------------------------------------

\subsection{Data Sources}

The forecasting pipelines combine ground-level pollutant observations with meteorological,
satellite-derived, and forecast-related variables. The data are aligned to an hourly time
grid before being converted into model inputs.

\begin{table}[htbp]
\centering
\caption{Main Data Sources Used by StratoWatch}
\label{tab:data_sources}
\begin{tabular}{p{0.25\columnwidth} p{0.34\columnwidth} p{0.29\columnwidth}}
\toprule
\textbf{Source} & \textbf{Variables} & \textbf{Purpose} \\
\midrule
Ground stations & Hourly O$_3$, NO$_2$ & Target observations \\
Meteorology & Temperature, humidity, wind, boundary-layer variables & Atmospheric context \\
Satellite data & Trace-gas observations & Regional chemical context \\
Forecast variables & Forecast-related pollutant/environmental variables & Additional predictive information \\
\bottomrule
\end{tabular}
\end{table}

Seven monitoring stations are used for the final multi-site experiment. The final single-site
pipeline uses 134 input features, while the multi-site representation contains 28 features per
station.

\subsection{Preprocessing}

The data are aligned to an hourly time grid and transformed into fixed-length historical
windows. Feature scaling is performed using the appropriate training-derived statistics.

The feature construction includes temporal information and meteorological variables, with
lagged and rolling information used where available. Wind direction is represented in a form
that can be used by the forecasting models rather than treating the circular direction value
as an ordinary scalar.

A chronological 70/15/15 train/validation/test split is used for the model experiments. The
test data are kept separate from training and model-selection decisions.

\subsection{Normalization and Evaluation Units}

An important distinction exists between the two final experiments. The final single-site
evaluation operates in near-normalized target units because the official single-site dataset
is already normalized and the corresponding target scaler has mean values close to zero and
unit scale. Its MAE and RMSE therefore should not be interpreted directly as
$\mu$g/m$^3$.

The final multi-site results are inverse-scaled into real concentration units using the final
target scaler. The target statistics are approximately 31.967~$\mu$g/m$^3$ mean and
34.720~$\mu$g/m$^3$ standard deviation for O$_3$, and 39.133~$\mu$g/m$^3$ mean and
29.909~$\mu$g/m$^3$ standard deviation for NO$_2$.

This distinction is important when comparing results across the two experiments.

\subsection{Graph Construction}

Two graph variants are evaluated in addition to the plain ST-Transformer.

The static graph represents fixed spatial relationships between stations using geographic
proximity. Each station is connected to its nearest neighbouring stations using
distance-based weights.

The dynamic wind graph retains candidate spatial relationships but changes their effective
strength using the prevailing wind information. The purpose is to test whether a graph that
changes with atmospheric transport conditions provides useful information beyond the
meteorological variables already supplied to the model.

% -------------------------------------------------------
\section{Proposed Methodology}
% -------------------------------------------------------

\subsection{Single-Site Models}

Five final single-site model families are compared: Forecast-only Baseline, Random Forest,
XGBoost, TCN, LSTM, and a temporal Transformer.

The Forecast-only Baseline provides a reference point without an additional learned forecasting
model. Random Forest provides a non-neural tree ensemble, while XGBoost provides gradient
boosted decision trees. The sequence models use the same general 24-hour forecasting setting
but learn temporal representations differently.

The final XGBoost implementation is a direct-target model. It uses twelve independent
XGBRegressor models, corresponding to each combination of six forecast horizons and two target
pollutants. The input contains 24 hours of 134 features, resulting in 3,216 input values per
sample.

\subsection{Temporal Models}

The LSTM processes the historical sequence recurrently, allowing information from earlier
hours to influence the forecast. The TCN uses temporal convolutions to capture dependencies
over different receptive-field sizes.

The temporal Transformer first projects the input features into a latent representation and
then applies self-attention across the historical time dimension. The final model configuration
uses a model dimension of 128, four attention heads, two temporal layers, dropout of 0.1,
SmoothL1Loss, and AdamW optimization.

\subsection{Multi-Site ST-Transformer}

The multi-site ST-Transformer processes the seven monitoring stations jointly. Each station
provides 28 features per hour over the 24-hour input window. Temporal dependencies are modeled
before allowing information from the different stations to interact.

The model predicts the complete six-hour horizon directly rather than generating one step and
feeding that prediction back into the model. This avoids the accumulation of autoregressive
errors across the six forecast steps.

\subsection{Graph-Enhanced Models}

The Static Graph ST-Transformer adds fixed spatial information between stations. The Dynamic
Wind Graph ST-Transformer extends this idea by changing graph connectivity strength according
to wind conditions.

The purpose of these models is not to assume that a graph must improve performance. Instead,
the experiments test whether explicit spatial structure provides useful information once
multi-station observations and meteorological variables are already available.

\begin{table}[htbp]
\centering
\caption{Main Model Configuration}
\label{tab:hyper}
\begin{tabular}{p{0.50\columnwidth} p{0.38\columnwidth}}
\toprule
\textbf{Parameter} & \textbf{Value} \\
\midrule
Input window & 24 hours \\
Forecast horizon & 6 hours \\
Multi-site stations & 7 \\
Single-site features & 134 \\
Multi-site features/site & 28 \\
Targets & O$_3$, NO$_2$ \\
Transformer dimension & 128 \\
Attention heads & 4 \\
Temporal layers & 2 \\
Dropout & 0.1 \\
Loss & SmoothL1Loss \\
Optimizer & AdamW \\
Early stopping & Patience 5 \\
Random seed for Transformer models & 42 \\
\bottomrule
\end{tabular}
\end{table}

% -------------------------------------------------------
\section{Experiments and Evaluation}
% -------------------------------------------------------

\subsection{Experimental Setup}

The final results are taken from frozen evaluation artifacts produced after the model
experiments were completed. The research figures are generated from these consolidated
results; no additional training is performed during figure generation.

The single-site test set contains 6,464 samples. The multi-site evaluation uses the frozen
seven-station configuration and reports results after inverse scaling to real concentration
units.

\subsection{Evaluation Metrics}

Three standard metrics are used:

\begin{equation}
\mathrm{MAE} =
\frac{1}{N}\sum_{i=1}^{N}|y_i-\hat{y}_i|
\end{equation}

\begin{equation}
\mathrm{RMSE} =
\sqrt{\frac{1}{N}\sum_{i=1}^{N}(y_i-\hat{y}_i)^2}
\end{equation}

and

\begin{equation}
R^2 =
1-\frac{\sum_i(y_i-\hat{y}_i)^2}
{\sum_i(y_i-\bar{y})^2}.
\end{equation}

MAE measures the average absolute forecasting error, RMSE gives greater weight to larger
errors, and $R^2$ measures the fraction of variance explained relative to the mean predictor.

% -------------------------------------------------------
\section{Results and Discussion}
% -------------------------------------------------------

\subsection{Single-Site Results}

Table~\ref{tab:single_site} shows the final single-site comparison. All values are reported
in the near-normalized evaluation units described in Section IV-C.

\begin{table*}[htbp]
\centering
\caption{Final Single-Site Model Performance in Near-Normalized Units}
\label{tab:single_site}
\begin{tabular}{lccc}
\toprule
\textbf{Model} & \textbf{MAE} & \textbf{RMSE} & \textbf{$R^2$} \\
\midrule
Forecast-only Baseline & 1.2384 & 1.5116 & -3.0296 \\
Random Forest & 0.4392 & 0.6526 & 0.2489 \\
\textbf{XGBoost} & \textbf{0.4039} & \textbf{0.5906} & \textbf{0.3848} \\
TCN & 0.5588 & 0.7427 & 0.0273 \\
LSTM & 0.6616 & 0.8296 & -0.2138 \\
Transformer & 0.5700 & 0.7524 & 0.0017 \\
\bottomrule
\end{tabular}
\end{table*}

XGBoost gives the strongest overall single-site result across all three reported metrics.
Its MAE is 0.4039, compared with 0.4392 for Random Forest. Its $R^2$ of 0.3848 also exceeds
the other evaluated models.

The result is useful because the single-site input contains a relatively large number of
engineered features compared with the length of the historical window. In this setting,
the tree-based models appear to make more effective use of the tabular representation than
the sequence models evaluated here.

The Forecast-only Baseline performs substantially worse, with an MAE of 1.2384 and a negative
$R^2$ of -3.0296. This provides a useful reference for the learned models, although the
near-normalized evaluation units mean that the magnitude should not be interpreted as a
physical concentration error.

\begin{figure}[htbp]
\centering
\includegraphics[width=\columnwidth]{01_single_site_mae.png}
\caption{Final single-site MAE comparison across the evaluated models.}
\label{fig:single_mae}
\end{figure}

\begin{figure}[htbp]
\centering
\includegraphics[width=\columnwidth]{03_single_site_r2.png}
\caption{Final single-site $R^2$ comparison across the evaluated models.}
\label{fig:single_r2}
\end{figure}

\subsection{Single-Site Target and Horizon Analysis}

For the final XGBoost model, the O$_3$ MAE is 0.3045 while the NO$_2$ MAE is 0.5034. The
corresponding $R^2$ values are 0.0609 and 0.3561, respectively. This indicates that the model
explains more of the variation in NO$_2$ than in O$_3$ under the current evaluation setup.

Performance also changes with forecast horizon. XGBoost obtains an MAE of 0.3497 at H+1,
which gradually increases to 0.4451 at H+6. Its $R^2$ decreases from 0.5306 at H+1 to
0.2673 at H+6.

\begin{table}[htbp]
\centering
\caption{Final XGBoost Horizon-Wise Performance}
\label{tab:xgb_horizon}
\begin{tabular}{lccc}
\toprule
\textbf{Horizon} & \textbf{MAE} & \textbf{RMSE} & \textbf{$R^2$} \\
\midrule
H+1 & 0.3497 & 0.5157 & 0.5306 \\
H+2 & 0.3731 & 0.5469 & 0.4723 \\
H+3 & 0.3968 & 0.5810 & 0.4046 \\
H+4 & 0.4225 & 0.6145 & 0.3342 \\
H+5 & 0.4364 & 0.6303 & 0.2998 \\
H+6 & 0.4451 & 0.6448 & 0.2673 \\
\bottomrule
\end{tabular}
\end{table}

The gradual increase in error with horizon is expected in a short-term forecasting problem:
as the prediction moves further away from the observed history, uncertainty increases. The
result also shows that the first forecast hour contains considerably more predictable
information than the later hours.

\begin{figure}[htbp]
\centering
\includegraphics[width=\columnwidth]{05_single_site_horizon_mae.png}
\caption{XGBoost MAE across the six forecast horizons in the single-site experiment.}
\label{fig:single_horizon}
\end{figure}

\subsection{Multi-Site Results}

The multi-site results are reported in real concentration units after inverse scaling.
Table~\ref{tab:multi_site} compares the plain ST-Transformer with the two graph-enhanced
variants.

\begin{table}[htbp]
\centering
\caption{Final Multi-Site Model Performance in $\mu$g/m$^3$}
\label{tab:multi_site}
\begin{tabular}{lccc}
\toprule
\textbf{Model} & \textbf{MAE} & \textbf{RMSE} & \textbf{$R^2$} \\
\midrule
ST-Transformer & \textbf{18.8349} & 27.9747 & 0.3437 \\
Static Graph ST & 19.0384 & 27.0781 & 0.3850 \\
Dynamic Wind Graph ST & 19.0464 & \textbf{26.9455} & \textbf{0.3911} \\
\bottomrule
\end{tabular}
\end{table}

The three models are relatively close in overall performance, but they do not produce the
same ranking for every metric. The plain ST-Transformer obtains the lowest MAE at
18.8349~$\mu$g/m$^3$. The Static Graph ST-Transformer improves on it in RMSE and $R^2$,
while the Dynamic Wind Graph ST-Transformer gives the best RMSE of 26.9455~$\mu$g/m$^3$ and
the highest $R^2$ of 0.3911.

This is different from simply saying that graph modeling either works or does not work.
Under the current configuration, graph modeling provides a modest improvement in metrics that
are more sensitive to larger errors and explained variance, while the plain ST-Transformer
retains the best average absolute error.

\begin{figure}[htbp]
\centering
\includegraphics[width=\columnwidth]{06_multi_site_mae.png}
\caption{Final multi-site MAE comparison.}
\label{fig:multi_mae}
\end{figure}

\begin{figure}[htbp]
\centering
\includegraphics[width=\columnwidth]{08_multi_site_r2.png}
\caption{Final multi-site $R^2$ comparison.}
\label{fig:multi_r2}
\end{figure}

\subsection{Multi-Site Target Analysis}

For the ST-Transformer, O$_3$ obtains an MAE of 19.0729~$\mu$g/m$^3$ and an $R^2$ of 0.4078,
while NO$_2$ obtains an MAE of 18.5969~$\mu$g/m$^3$ and an $R^2$ of 0.2078.

For the Dynamic Wind Graph ST-Transformer, O$_3$ obtains an MAE of
18.7858~$\mu$g/m$^3$ and an $R^2$ of 0.4769. NO$_2$ obtains an MAE of
19.3071~$\mu$g/m$^3$ and an $R^2$ of 0.2226.

The difference between the pollutants is therefore not uniform across all metrics. The
dynamic graph performs particularly well for O$_3$ in terms of explained variance, while its
NO$_2$ MAE is slightly higher.

\begin{figure}[htbp]
\centering
\includegraphics[width=\columnwidth]{09_multi_site_o3_no2_mae.png}
\caption{O$_3$ and NO$_2$ MAE comparison for the final multi-site models.}
\label{fig:multi_target}
\end{figure}

\subsection{Multi-Site Horizon Analysis}

The Static Graph and Dynamic Wind Graph models show a gradual increase in MAE as the forecast
moves from H+1 to H+6. For the Dynamic Wind Graph ST-Transformer, MAE increases from
18.8561~$\mu$g/m$^3$ at H+1 to 19.2365~$\mu$g/m$^3$ at H+6.

Its $R^2$ remains relatively stable over the six horizons, ranging from 0.3975 at H+1 to
0.3752 at H+6. This suggests that although the absolute error increases with horizon, the
overall explained variance does not collapse over the six-hour prediction window.

\begin{figure}[htbp]
\centering
\includegraphics[width=\columnwidth]{10_multi_site_horizon_mae.png}
\caption{Horizon-wise MAE for the final multi-site graph models.}
\label{fig:multi_horizon_mae}
\end{figure}

% -------------------------------------------------------
\section{Qualitative Analysis}
% -------------------------------------------------------

Quantitative metrics describe average behaviour, but they do not show exactly where a model
struggles. The final evaluation therefore includes prediction-versus-actual plots for the
pollutants.

Figure~\ref{fig:prediction_o3} shows the final O$_3$ prediction comparison, while
Figure~\ref{fig:prediction_no2} shows the corresponding NO$_2$ result.

\begin{figure}[htbp]
\centering
\includegraphics[width=\columnwidth]{12_prediction_vs_actual_o3.png}
\caption{Prediction versus actual O$_3$ values from the final evaluation.}
\label{fig:prediction_o3}
\end{figure}

\begin{figure}[htbp]
\centering
\includegraphics[width=\columnwidth]{12_prediction_vs_actual_no2.png}
\caption{Prediction versus actual NO$_2$ values from the final evaluation.}
\label{fig:prediction_no2}
\end{figure}

The final evaluation also shows a tendency to underestimate some high-pollution episodes.
This behaviour matters because an early-warning system is particularly sensitive to missed
peaks. Average error metrics alone can hide this failure mode, especially when most samples
are closer to normal concentration levels.

The present results therefore suggest that future work should evaluate models not only by
overall MAE and RMSE, but also by their behaviour during high-concentration events.

% -------------------------------------------------------
\section{System Implementation}
% -------------------------------------------------------

The research models are integrated into the StratoWatch web application. The frontend is
implemented using React and Vite, while the backend uses FastAPI. The backend exposes separate
interfaces for the single-site and multi-site forecasting pipelines.

The application allows a user to provide input data, select the forecasting workflow, and
inspect the resulting predictions, metrics, and visualizations. The final backend uses the
frozen model artifacts rather than starting a new training process when a forecast request is
made.

This separation is useful for the research workflow because the models used to produce the
reported results can remain fixed while the web interface is used for inference and
demonstration.

The deployed system is therefore not intended to replace the research evaluation pipeline.
Instead, it provides a practical interface around the finalized forecasting models and makes
the experimental work easier to inspect and demonstrate.

% -------------------------------------------------------
\section{Limitations and Future Work}
% -------------------------------------------------------

The most important limitation of the current study is the six-hour forecast horizon. The
problem motivating the work targets substantially longer operational forecasts, including
24--48 hours. The current results should therefore not be interpreted as evidence that the
system already satisfies that operational requirement.

A second limitation is the difference in evaluation units between the two experiments.
Single-site metrics are reported in near-normalized units, while the multi-site metrics are
reported in real concentration units. This prevents direct numerical comparison of the MAE
values between the two experiments.

The multi-site study is also limited to seven monitoring stations. The final graph checkpoints
were trained and evaluated for this configuration, so conclusions about other station layouts
or larger geographic areas require additional experiments.

Another limitation is the difficulty of forecasting high-pollution spikes. The current models
can underestimate extreme values, which is particularly important for an early-warning
application. The present study does not include a dedicated extreme-event objective or
uncertainty calibration.

Finally, the dynamic graph is based on available wind information and does not model the full
vertical structure of atmospheric transport. Real pollutant movement can depend on processes
that are not captured by a surface-wind-based relationship alone.

Future work will focus on:

\begin{enumerate}
\item Extending the forecasting horizon to 24 and 48 hours.
\item Improving the prediction of high-pollution episodes using spike-aware objectives or
sampling strategies.
\item Adding calibrated predictive uncertainty to the forecasts.
\item Learning spatial relationships directly from data rather than relying only on
predefined geographic and wind-based relationships.
\item Incorporating additional emission-related information such as traffic, fire, and
industrial activity proxies where suitable data are available.
\item Connecting the deployed system to continuously updated meteorological and
satellite feeds for genuinely real-time forecasting.
\end{enumerate}

% -------------------------------------------------------
\section{Conclusion}
% -------------------------------------------------------

StratoWatch presents a practical study of short-horizon O$_3$ and NO$_2$ forecasting at both
single-site and multi-site scales. At the single-site level, XGBoost gives the strongest
overall result among the evaluated models, with an MAE of 0.4039, RMSE of 0.5906, and
$R^2$ of 0.3848 in the near-normalized evaluation units.

The multi-site experiment gives a more mixed picture. The plain ST-Transformer obtains the
lowest MAE at 18.8349~$\mu$g/m$^3$, while the Dynamic Wind Graph ST-Transformer obtains the
lowest RMSE and highest $R^2$. The Static Graph ST-Transformer remains competitive. These
results suggest that explicit spatial relationships can provide useful information, but the
benefit is metric-dependent rather than a uniform improvement over the non-graph model.

The experiments also show that forecasting becomes more difficult as the prediction horizon
increases and that high-pollution episodes remain challenging. These observations are
important because a forecasting system intended for early warning should eventually be
evaluated specifically on the events that matter most operationally.

The current StratoWatch system is therefore best viewed as a short-horizon research prototype
and evaluation platform. Its main value at this stage is not claiming a finished operational
forecasting solution, but providing a reproducible comparison of several forecasting
approaches and a working system in which those models can be evaluated and demonstrated.

% -------------------------------------------------------
\section*{Data and Code Availability}
% -------------------------------------------------------

The StratoWatch source code, model implementations, preprocessing pipeline, evaluation
artifacts, and web application are maintained in the project's public GitHub repository.

Ground-level air-quality observations were obtained from CPCB-related sources, while
satellite-derived atmospheric information is based on Sentinel-5P/TROPOMI products and other
environmental data used by the project. Redistribution of third-party data remains subject
to the terms and conditions of the respective providers.

% -------------------------------------------------------
\section*{Acknowledgment}
% -------------------------------------------------------

The author thanks the open-source communities behind PyTorch, XGBoost, scikit-learn,
FastAPI, React, and the broader scientific software ecosystem used in this work. The author
also acknowledges the providers of the ground-level air-quality, meteorological, and
satellite-derived datasets used in the study.

% -------------------------------------------------------
% REFERENCES
% -------------------------------------------------------

\begin{thebibliography}{00}

\bibitem{vaswani}
A. Vaswani, N. Shazeer, N. Parmar, J. Uszkoreit, L. Jones, A. N. Gomez,
L. Kaiser, and I. Polosukhin, ``Attention is all you need,'' in
\emph{Advances in Neural Information Processing Systems (NeurIPS)}, vol. 30,
2017, pp. 5998--6008.

\bibitem{xgboost}
T. Chen and C. Guestrin, ``XGBoost: A scalable tree boosting system,'' in
\emph{Proc. 22nd ACM SIGKDD International Conference on Knowledge Discovery
and Data Mining (KDD)}, 2016, pp. 785--794.

\bibitem{lstm}
S. Hochreiter and J. Schmidhuber, ``Long short-term memory,''
\emph{Neural Computation}, vol. 9, no. 8, pp. 1735--1780, Nov. 1997.

\bibitem{tcn}
S. Bai, J. Z. Kolter, and V. Koltun, ``An empirical evaluation of generic
convolutional and recurrent networks for sequence modeling,''
\emph{arXiv preprint arXiv:1803.01271}, 2018.

\bibitem{graphwavenet}
Z. Wu, S. Pan, G. Long, J. Jiang, and C. Zhang, ``Graph WaveNet for deep
spatial-temporal graph modeling,'' in \emph{Proc. 28th International Joint
Conference on Artificial Intelligence (IJCAI)}, 2019, pp. 1907--1913.

\bibitem{informer}
H. Zhou, S. Zhang, J. Peng, J. Zhang, J. Li, W. Xiong, and W. Zhang,
``Informer: Beyond efficient Transformer for long sequence time-series
forecasting,'' in \emph{Proc. AAAI Conference on Artificial Intelligence},
vol. 35, no. 12, 2021, pp. 11106--11115.

\bibitem{agcrn}
L. Bai, L. Yao, C. Li, X. Wang, and C. Wang, ``Adaptive graph convolutional
recurrent network for traffic forecasting,'' in
\emph{Advances in Neural Information Processing Systems (NeurIPS)},
vol. 33, 2020.

\bibitem{dcrnn}
Y. Li, R. Yu, C. Shahabi, and Y. Liu, ``Diffusion convolutional recurrent
neural network: Data-driven traffic forecasting,'' in
\emph{Proc. International Conference on Learning Representations (ICLR)}, 2018.

\bibitem{ceemdan2025}
S. He, Y. Liu, J. Peng, D. Xu, and Y. Wang, ``A CEEMDAN-GNN-transformer
hybrid model for air quality index forecasting: a case study of Chang'an
town, Dongguan, China,'' \emph{Frontiers in Environmental Science},
vol. 13, art. 1653446, 2025.

\bibitem{estgcn2025}
M. Panja, T. Chakraborty, A. Biswas, and S. Deb, ``E-STGCN: Extreme
spatiotemporal graph convolutional networks for air quality forecasting,''
\emph{arXiv preprint arXiv:2411.12258}, 2025.

\end{thebibliography}

\end{document}